\chapter{Appendix: Summations}

You have already met summations everywhere in this book. Every time you wrote
``the loop runs once per element, and the $i$-th pass costs about $i$ steps,''
and then asked ``so how much in total?'', you were really asking for a
\emph{summation}. The analysis of almost every algorithm bottoms out in adding a
list of numbers together and asking how big the answer is.

This appendix is a toolbox. It collects the handful of sums that show up again
and again --- arithmetic, geometric, harmonic --- proves the formulas you have
been quoting, and then teaches you the three reliable \emph{techniques} for
sums you do not have a formula for: telescoping, splitting, and bounding by an
integral. By the end you should be able to look at any of the recurrences in
this book, expand it into a sum, and squeeze a clean $\Theta$ out of it.

\begin{focusbox}
A summation is just repeated addition with a name. The skill is not adding ---
it is recognizing the \emph{shape} of a sum so you can replace a loop's worth of
work with a single closed-form expression or a tight bound.
\end{focusbox}

\section{Notation: what \texorpdfstring{$\sum$}{the sigma} actually says}
\label{ch37:notation}

Let us be very explicit, because sloppy reading of $\sum$ causes more grief than
the arithmetic itself. The expression
\[
  \sum_{k=1}^{n} a_k \;=\; a_1 + a_2 + a_3 + \cdots + a_n
\]
means: ``start the counter $k$ at $1$; for each integer value of $k$ up to and
including $n$, compute the \emph{term} $a_k$; add all those terms together.'' The
letter $k$ is the \emph{index of summation}, $1$ is the \emph{lower limit}, $n$
is the \emph{upper limit}, and $a_k$ is the term, sometimes called the
\emph{summand}.

Three things trip people up, so pin them down now.

\begin{itemize}
  \item \textbf{The index is a dummy.} $\sum_{k=1}^{n} a_k$ and
        $\sum_{i=1}^{n} a_i$ and $\sum_{j=1}^{n} a_j$ are the \emph{same number}.
        The index name vanishes after you add; it never escapes the sum.
  \item \textbf{An empty sum is $0$.} If the upper limit is below the lower limit,
        for example $\sum_{k=1}^{0} a_k$, there are no terms, and the convention is
        that the total is $0$. (Likewise an empty \emph{product} is $1$.) This is
        not a special case to fear; it is the choice that makes formulas keep
        working at the boundary.
  \item \textbf{Count the terms.} The sum $\sum_{k=m}^{n}$ has exactly
        $n - m + 1$ terms, \emph{not} $n - m$. From $k=3$ to $k=7$ inclusive is
        $7 - 3 + 1 = 5$ terms: $3,4,5,6,7$. Off-by-one errors in term-counting are
        the number-one source of wrong constants.
\end{itemize}

\begin{trapbox}
``$\sum_{k=1}^{n} 1 = 1$.'' No. The summand happens to be the constant $1$, but
you still add it once per value of $k$, so $\sum_{k=1}^{n} 1 = n$. A constant
summand $c$ gives $\sum_{k=1}^{n} c = c\,n$. Always ask ``how many terms?''
before you ask ``what is each term?''.
\end{trapbox}

\subsection{Two properties you will use constantly}
\label{ch37:linearity}

Summation is \emph{linear}, which is a fancy way of saying you can pull constants
out and split sums apart. For any constant $c$ and any terms $a_k, b_k$,
\[
  \sum_{k=1}^{n} (c\,a_k + b_k)
    \;=\; c\sum_{k=1}^{n} a_k \;+\; \sum_{k=1}^{n} b_k .
\]
That is the whole rule. Both halves follow from the commutativity and
associativity of ordinary addition --- you are allowed to reorder and regroup a
finite list of numbers however you like. Linearity is what lets you turn a messy
summand into a combination of sums you already know.

A second move, \emph{re-indexing} (also called shifting), changes which value the
counter starts at without changing the total. If you substitute $j = k - 1$, then
when $k$ runs from $1$ to $n$, the new index $j$ runs from $0$ to $n-1$:
\[
  \sum_{k=1}^{n} a_k \;=\; \sum_{j=0}^{n-1} a_{j+1}.
\]
Same terms, relabelled. Re-indexing is the trick that makes telescoping and
geometric-series manipulations work; we will lean on it below.

\begin{example}[Linearity in action]
\label{ch37:ex-linearity}
Suppose you need $\sum_{k=1}^{n}(3k + 2)$. Do not panic and do not expand by
hand. Split it:
\[
  \sum_{k=1}^{n}(3k+2) = 3\sum_{k=1}^{n} k + \sum_{k=1}^{n} 2
                       = 3\cdot\frac{n(n+1)}{2} + 2n,
\]
using the arithmetic-series formula we prove next. For $n=4$ this is
$3\cdot 10 + 8 = 38$, and indeed $5+8+11+14 = 38$. The structure did the work.
\end{example}

\section{The arithmetic series}
\label{ch37:arithmetic}

The most fundamental sum in algorithm analysis is $1 + 2 + 3 + \cdots + n$. It is
what you get when a loop's inner work grows by one each pass --- selection sort,
insertion sort's worst case, the nested loops of a naive comparison routine. Its
closed form is the one formula in this appendix you should be able to write in
your sleep.

\begin{intuitionbox}
Picture the numbers $1$ through $n$ written left to right, then the \emph{same}
numbers written right to left underneath:
\[
\begin{array}{cccccc}
 1 & 2 & 3 & \cdots & n-1 & n \\
 n & n-1 & n-2 & \cdots & 2 & 1
\end{array}
\]
Each \emph{column} adds to $n+1$. There are $n$ columns, so both rows together
total $n(n+1)$. But the two rows are the same sum, so one row is half of that.
That is the legendary ``fold the list in half'' trick, attributed to a young
Gauss, and it \emph{is} the proof.
\end{intuitionbox}

\begin{theorem}[Sum of the first $n$ integers]
\label{ch37:thm-arith}
For every integer $n \ge 0$,
\[
  \sum_{k=1}^{n} k \;=\; \frac{n(n+1)}{2}.
\]
\end{theorem}

\begin{proof}
We argue by induction on $n$, which mirrors the way a loop builds the sum one
term at a time. \emph{Base case:} for $n=0$ the sum is empty, equal to $0$, and
the formula gives $\tfrac{0\cdot 1}{2}=0$. They agree.

\emph{Inductive step:} assume $\sum_{k=1}^{n} k = \tfrac{n(n+1)}{2}$ for some
$n \ge 0$. Then add the next term $n+1$:
\[
  \sum_{k=1}^{n+1} k
    = \Bigl(\sum_{k=1}^{n} k\Bigr) + (n+1)
    = \frac{n(n+1)}{2} + (n+1)
    = (n+1)\cdot\frac{n + 2}{2}
    = \frac{(n+1)\bigl((n+1)+1\bigr)}{2}.
\]
That is exactly the formula with $n$ replaced by $n+1$, so by induction it holds
for all $n \ge 0$.
\end{proof}

The same fold-in-half idea works for \emph{any} arithmetic series --- one whose
terms step up by a constant difference $d$. If the first term is $a$ and there are
$m$ terms, the last term is $a + (m-1)d$, and
\[
  \sum_{k=0}^{m-1}(a + kd)
    \;=\; m\cdot\underbrace{\frac{(\text{first}) + (\text{last})}{2}}_{\text{average term}}
    \;=\; \frac{m\bigl(2a + (m-1)d\bigr)}{2}.
\]
Read it as ``number of terms times the average of the first and last.'' That
sentence is worth more than the formula, because it survives even when you have
forgotten the algebra.

\begin{example}[Tracing a nested loop]
\label{ch37:ex-nested}
Consider this fragment, the heart of selection sort:
\begin{verbatim}
for i = 1 to n:
    for j = i to n:
        do one unit of work
\end{verbatim}
The inner loop runs $n - i + 1$ times when the outer index is $i$. Total work:
\[
  \sum_{i=1}^{n}(n - i + 1)
    = \sum_{t=1}^{n} t
    = \frac{n(n+1)}{2},
\]
where we re-indexed with $t = n - i + 1$ (as $i$ goes $1\to n$, $t$ goes
$n\to 1$, the same set of values). So the running time is
$\Theta(n^2)$ --- and now you know the leading constant is $\tfrac12$, not $1$.
\end{example}

\begin{aestheticbox}
Why care about the $\tfrac12$? Because constants decide real races. Two
$\Theta(n^2)$ sorts can differ by $2\times$ in wall-clock time purely from
whether the inner loop is $\sum i$ or $\sum 2i$. Asymptotics tell you which
algorithms to throw away; summation constants tell you which of the survivors to
ship.
\end{aestheticbox}

\subsection{Sums of squares and higher powers}
\label{ch37:squares}

One step up in difficulty is the sum of squares, which appears whenever the inner
cost itself grows quadratically (think of summing the comparison counts in some
divide-and-combine routines):
\[
  \sum_{k=1}^{n} k^2 \;=\; \frac{n(n+1)(2n+1)}{6}.
\]
You can prove it by the same induction; the bookkeeping is heavier but the idea is
identical. The pattern continues: $\sum_{k=1}^{n} k^p$ is a polynomial in $n$ of
degree $p+1$ with leading term $\tfrac{n^{p+1}}{p+1}$. That single fact ---
\[
  \sum_{k=1}^{n} k^p \;=\; \Theta\!\bigl(n^{p+1}\bigr)\quad\text{for any fixed } p \ge 0
\]
--- is usually all an algorithm analysis needs, and in
Section~\ref{ch37:integrals} we will see \emph{why} the $\tfrac{1}{p+1}$ leading
constant has to be there, using nothing but the area under a curve.

\section{The geometric series}
\label{ch37:geometric}

The second pillar is the geometric series, where each term is a fixed ratio $x$
times the one before: $1, x, x^2, x^3, \dots$. This is the sum behind binary
representations, doubling arrays, balanced-tree level counts, and the
``merge'' levels of merge sort. It behaves \emph{completely differently}
depending on whether $x<1$ or $x>1$, and understanding that split is the single
most useful thing in this appendix.

\begin{theorem}[Finite geometric series]
\label{ch37:thm-geo}
For any real $x \neq 1$ and integer $n \ge 0$,
\[
  \sum_{k=0}^{n} x^k \;=\; 1 + x + x^2 + \cdots + x^n \;=\; \frac{x^{n+1} - 1}{x - 1}.
\]
\end{theorem}

\begin{proof}
This is the prettiest \emph{telescoping} argument in the book, so savour it. Call
the sum $S = \sum_{k=0}^{n} x^k$. Multiply both sides by $x$ and re-index:
\[
  xS = \sum_{k=0}^{n} x^{k+1} = \sum_{j=1}^{n+1} x^{j} = x + x^2 + \cdots + x^{n+1}.
\]
Now subtract. Almost every term cancels --- the $x, x^2, \dots, x^n$ appear in
both $S$ and $xS$:
\[
  xS - S = x^{n+1} - x^0 = x^{n+1} - 1,
  \qquad\text{so}\qquad
  S(x-1) = x^{n+1} - 1.
\]
Divide by $x - 1$ (legal because $x \neq 1$) to finish. The trick was to create a
\emph{shifted copy} of the sum so that subtracting collapses it.
\end{proof}

\begin{focusbox}
Memorize the proof move, not just the answer: \emph{multiply by the ratio, shift,
subtract, watch the middle cancel.} You will reuse it for sums like
$\sum k\,x^k$ that have no formula on the tip of your tongue.
\end{focusbox}

\subsection{The two regimes: shrinking vs.\ growing}
\label{ch37:geo-regimes}

Everything interesting about geometric series lives in one question: is the ratio
$x$ less than $1$ or greater than $1$?

\textbf{Shrinking ($0 < x < 1$).} The terms get smaller and the sum
\emph{converges} even if you let $n \to \infty$. Take the limit of the formula:
since $x^{n+1} \to 0$,
\[
  \sum_{k=0}^{\infty} x^k \;=\; \frac{1}{1-x}.
\]
The headline consequence for algorithms: a shrinking geometric sum is
\emph{dominated by its first term}. For instance
$\sum_{k=0}^{\infty} (1/2)^k = 2$, and the very first term is $1$ --- the entire
infinite tail only adds another $1$. So when costs halve at each level, the
\emph{top} level pays for everything, up to a constant factor.

\textbf{Growing ($x > 1$).} Now the terms blow up and the sum is
\emph{dominated by its last term}. From the formula,
\[
  \sum_{k=0}^{n} x^k = \frac{x^{n+1}-1}{x-1} \le \frac{x^{n+1}}{x-1}
                     = x^{n}\cdot\frac{x}{x-1} = \Theta\bigl(x^{n}\bigr).
\]
The last term $x^n$ already is $\Theta(x^n)$, and everything below it adds only a
constant multiple. So when costs double at each level, the \emph{bottom} level
pays for everything, up to a constant factor.

\begin{center}
\begin{tikzpicture}[scale=1.0]
  % --- shrinking regime ---
  \begin{scope}
    \node[font=\scriptsize\bfseries] at (1.75,2.55) {shrinking $x=\tfrac12$};
    \foreach \k/\h in {0/2.0,1/1.0,2/0.5,3/0.25,4/0.125}{
      \draw[fill=sourceblue!25,draw=sourceblue!70!black]
        (\k*0.7,0) rectangle (\k*0.7+0.55,\h);
    }
    \draw[->] (-0.15,0) -- (3.7,0);
    \node[font=\scriptsize,align=center] at (1.75,-0.45)
      {first term (the tallest) dominates};
  \end{scope}
  % --- growing regime ---
  \begin{scope}[xshift=5.2cm]
    \node[font=\scriptsize\bfseries] at (1.75,2.55) {growing $x=2$};
    \foreach \k/\h in {0/0.125,1/0.25,2/0.5,3/1.0,4/2.0}{
      \draw[fill=quotientorange!30,draw=quotientorange!75!black]
        (\k*0.7,0) rectangle (\k*0.7+0.55,\h);
    }
    \draw[->] (-0.15,0) -- (3.7,0);
    \node[font=\scriptsize,align=center] at (1.75,-0.45)
      {last term (the tallest) dominates};
  \end{scope}
\end{tikzpicture}
\end{center}

Each bar is one term $x^k$, drawn left to right as $k$ increases. When $x<1$ the
bars shrink and the leftmost (first) bar towers over the rest; when $x>1$ they
grow and the rightmost (last) bar does. In both cases the tallest bar alone is a
constant fraction of the whole sum --- that is exactly what ``dominated by one
term, up to a constant factor'' means.

\begin{intuitionbox}
This ``first term wins'' / ``last term wins'' dichotomy is the secret engine of
the Master Theorem. A recursion tree has geometric per-level costs; whether the
root or the leaves dominate is \emph{exactly} the question of whether the ratio is
below or above $1$. If the per-level cost shrinks, the root wins (case akin to
$\Theta(f(n))$); if it grows, the leaves win (case akin to $\Theta(n^{\log_b a})$);
if it stays flat ($x = 1$), every level costs the same and you pay an extra factor
of the number of levels, $\Theta(\log n)$.
\end{intuitionbox}

\begin{example}[Building a dynamic array]
\label{ch37:ex-doubling}
A growable array that \emph{doubles} its capacity copies $1, 2, 4, \dots, n$
elements over its lifetime as it grows to size $n$. The total copy work is a
\emph{growing} geometric series:
\[
  1 + 2 + 4 + \cdots + n = \sum_{k=0}^{\log_2 n} 2^k = 2^{\log_2 n + 1} - 1 = 2n - 1.
\]
Last term wins: the grand total is just under $2n$, even though you resized
$\log_2 n$ times. Spread over $n$ insertions, that is $O(1)$ \emph{amortized}
cost per insertion. The doubling strategy is cheap precisely because the
geometric series is dominated by its final, biggest copy.
\end{example}

\section{The harmonic series}
\label{ch37:harmonic}

The third pillar is subtler and, frankly, surprising the first time you meet it.
The \emph{harmonic series} is
\[
  H_n \;=\; \sum_{k=1}^{n} \frac{1}{k} \;=\; 1 + \frac12 + \frac13 + \cdots + \frac1n .
\]
Each term shrinks toward $0$, so you might guess --- as with a shrinking geometric
series --- that the total settles to a finite limit. It does not. $H_n$ grows
without bound, just \emph{very} slowly, like $\ln n$.

\begin{intuitionbox}
Group the terms in blocks whose sizes double: $\{1\}$, $\{\tfrac12\}$,
$\{\tfrac13,\tfrac14\}$, $\{\tfrac15,\dots,\tfrac18\}$, and so on. In the block
ending at $\tfrac{1}{2^m}$ there are $2^{m-1}$ terms, each at least
$\tfrac{1}{2^m}$, so the block sums to at least $2^{m-1}\cdot\tfrac{1}{2^m} =
\tfrac12$. Every doubling block contributes at least another $\tfrac12$. To get a
total of $T$ you need about $2T$ doubling blocks, i.e.\ about $2^{2T}$ terms. The
sum grows, but you must go exponentially far to gain a constant --- that is
logarithmic growth, seen from the inside.
\end{intuitionbox}

\begin{theorem}[Harmonic growth]
\label{ch37:thm-harmonic}
$H_n = \ln n + \Theta(1)$. More precisely, $\ln n < H_n \le \ln n + 1$, and in
fact $H_n = \ln n + \gamma + o(1)$ where $\gamma \approx 0.5772$ is the
Euler--Mascheroni constant. For algorithm analysis the headline is simply
\[
  H_n = \Theta(\log n).
\]
\end{theorem}

We will prove the $\ln n < H_n \le \ln n + 1$ bounds cleanly in the next section
using integrals --- it is the cleanest possible illustration of the
integral-bounding technique.

\begin{aestheticbox}
The harmonic series is the running time of the \emph{average case} of quicksort
and of building a randomized binary search tree, and it governs the ``coupon
collector'' problem (how many random draws to see all $n$ coupons: about
$n H_n \approx n\ln n$). Whenever a cost looks like ``$1 + \tfrac12 + \tfrac13 +
\cdots$,'' you should immediately think $\Theta(\log n)$ and move on.
\end{aestheticbox}

\section{Bounding a sum by an integral}
\label{ch37:integrals}

Here is the most \emph{general} tool in the appendix, the one that rescues you
when no closed form exists. The idea is geometric and you should carry the picture
in your head forever: a sum $\sum f(k)$ is a row of rectangles, each of width $1$
and height $f(k)$, and the integral $\int f(x)\,dx$ is the smooth area under the
curve $f$. If you nest the rectangles against the curve correctly, the integral
\emph{traps} the sum from above and below.

\begin{center}
\begin{tikzpicture}[scale=1.0]
  % axes
  \draw[->] (-0.2,0) -- (7.0,0) node[right] {$x$};
  \draw[->] (0,-0.2) -- (0,3.4) node[above] {$y$};
  % a decreasing curve f(x) = 2.6/sqrt(x)
  \draw[thick,sourceblue,domain=0.55:6.5,samples=100,smooth]
        plot (\x, {2.6/sqrt(\x)}) node[right,sourceblue] {$y=f(x)$};
  % lower rectangles: rectangle on [k, k+1] has height f(k), so it sits
  % UNDER the decreasing curve. These rectangles total sum_{k} f(k) and
  % under-estimate the area, giving sum f(k) <= integral.
  \foreach \k in {1,2,3,4,5}{
     \pgfmathsetmacro{\h}{2.6/sqrt(\k)}
     \draw[fill=green!12,draw=green!55!black]
        (\k,0) rectangle (\k+1,\h);
  }
  % tick marks and index labels on the x-axis
  \foreach \k in {1,2,3,4,5,6}{
     \draw (\k,0.06) -- (\k,-0.06);
     \node[below] at (\k,-0.05) {\scriptsize $\k$};
  }
  % label one representative rectangle's height
  \draw[<-,gray!70] (1.5,1.4) -- (1.95,2.0)
        node[right,black] {\scriptsize height $f(k)$, width $1$};
  \node[align=center] at (4.4,2.9)
     {\scriptsize each rectangle has area $f(k)$, so the staircase\\[-1pt]
      \scriptsize total $\sum f(k)$ sits just below $\int f(x)\,dx$};
\end{tikzpicture}
\end{center}

The picture above shows the easy half for a decreasing $f$: each rectangle of
height $f(k)$ tucks completely \emph{under} the curve, so the staircase area
$\sum f(k)$ can never exceed the smooth area $\int f(x)\,dx$ that covers it. Slide
the same rectangles one step the other way and they poke \emph{above} the curve,
giving the matching bound from the other side. That single picture --- rectangles
trapped between two slices of the curve --- is the whole technique.

The rule has two halves, one for increasing functions and one for decreasing
functions. State it once and you never have to rederive it.

\begin{formalbox}
\textbf{Integral bounds.} Let $f$ be \emph{monotone} on the relevant range.
\begin{itemize}
  \item If $f$ is \emph{increasing}, then
        \[
          \int_{m-1}^{n} f(x)\,dx \;\le\; \sum_{k=m}^{n} f(k) \;\le\; \int_{m}^{n+1} f(x)\,dx.
        \]
  \item If $f$ is \emph{decreasing}, then
        \[
          \int_{m}^{n+1} f(x)\,dx \;\le\; \sum_{k=m}^{n} f(k) \;\le\; \int_{m-1}^{n} f(x)\,dx.
        \]
\end{itemize}
The mantra: \emph{the rectangle of height $f(k)$ sits between the lower and upper
slices of curve next to it.} You rarely need to remember which inequality is
which --- just draw the rectangles against the curve and read off the picture.
\end{formalbox}

\begin{example}[Why $\sum k^p = \Theta(n^{p+1})$]
\label{ch37:ex-poly-integral}
Let $f(x) = x^p$ for a fixed $p \ge 0$; it is increasing, so
\[
  \int_{0}^{n} x^p\,dx \;\le\; \sum_{k=1}^{n} k^p \;\le\; \int_{1}^{n+1} x^p\,dx.
\]
The left integral is $\tfrac{n^{p+1}}{p+1}$ and the right is
$\tfrac{(n+1)^{p+1}-1}{p+1}$. Both are $\Theta(n^{p+1})$ with leading constant
$\tfrac{1}{p+1}$, so
\[
  \sum_{k=1}^{n} k^p = \frac{n^{p+1}}{p+1} + \Theta(n^p) = \Theta\bigl(n^{p+1}\bigr).
\]
No clever induction, no guessing the polynomial --- the area under $x^p$
\emph{is} the answer.
\end{example}

\begin{example}[Pinning down the harmonic series]
\label{ch37:ex-harmonic-integral}
Now take $f(x) = 1/x$, which is decreasing. The decreasing rule gives
\[
  \int_{1}^{n+1}\frac{dx}{x} \;\le\; \sum_{k=1}^{n}\frac1k \;\le\; 1 + \int_{1}^{n}\frac{dx}{x}.
\]
(The extra ``$1$'' on the right is the first term $f(1)=1$, peeled off so the
remaining terms line up under the decreasing curve from $1$ to $n$.) Since
$\int_1^{t} \tfrac{dx}{x} = \ln t$, we get
\[
  \ln(n+1) \;\le\; H_n \;\le\; 1 + \ln n.
\]
So $H_n$ is squeezed between $\ln n$ and $\ln n + 1$: it is $\ln n$ plus a number
that never escapes the interval $[0,1]$. That is the promised proof of
Theorem~\ref{ch37:thm-harmonic}, and it is genuinely just ``the staircase hugs
the logarithm.''
\end{example}

\begin{walkbox}
Take $f(x) = \ln x$ on a walk. Use the integral bounds to show that
$\sum_{k=1}^{n} \ln k$ --- which is $\ln(n!)$ --- is $\Theta(n \ln n)$. You have
just derived the asymptotic heart of \emph{Stirling's formula}, the reason
comparison sorting cannot beat $n\log n$. (Hint: $\int \ln x\,dx = x\ln x - x$.)
\end{walkbox}

\section{Telescoping: when the middle cancels}
\label{ch37:telescoping}

A \emph{telescoping} sum is one whose terms are built as a difference
$a_{k} - a_{k-1}$ (or $a_{k+1} - a_k$), so that consecutive pieces annihilate and
only the two endpoints survive --- like a collapsing spyglass. Formally,
\[
  \sum_{k=1}^{n}\bigl(a_k - a_{k-1}\bigr) \;=\; a_n - a_0 .
\]
You already saw a disguised telescope: the geometric-series proof subtracted $S$
from $xS$ and watched the middle collapse. Telescoping is your friend whenever the
summand looks like a difference, and your goal whenever you can engineer one.

\begin{example}[A partial-fractions telescope]
\label{ch37:ex-telescope}
Evaluate $\displaystyle\sum_{k=1}^{n}\frac{1}{k(k+1)}$. The key observation is the
partial-fraction split
\[
  \frac{1}{k(k+1)} = \frac{1}{k} - \frac{1}{k+1}.
\]
So the sum telescopes:
\[
  \sum_{k=1}^{n}\!\left(\frac1k - \frac{1}{k+1}\right)
    = \left(1 - \tfrac12\right) + \left(\tfrac12 - \tfrac13\right) + \cdots
      + \left(\tfrac1n - \tfrac{1}{n+1}\right)
    = 1 - \frac{1}{n+1}.
\]
Every interior fraction is created once with a $+$ and destroyed once with a $-$.
The diagram below shows the bookkeeping: each $-\tfrac1{k+1}$ from one term
pairs off against the $+\tfrac1{k+1}$ from the next, leaving only the very first
$+1$ and the very last $-\tfrac{1}{n+1}$ unmatched.

\begin{center}
\begin{tikzpicture}[scale=1.0, every node/.style={font=\scriptsize}]
  % the chain of terms (1 - 1/2) + (1/2 - 1/3) + ...
  \node (a1) at (0.0,0)   {$+1$};
  \node (b1) at (0.9,0)   {$-\tfrac12$};
  \node (a2) at (2.0,0)   {$+\tfrac12$};
  \node (b2) at (2.9,0)   {$-\tfrac13$};
  \node (a3) at (4.0,0)   {$+\tfrac13$};
  \node (b3) at (4.9,0)   {$-\tfrac14$};
  \node (d)  at (5.9,0)   {$\cdots$};
  \node (an) at (7.1,0)   {$+\tfrac1n$};
  \node (bn) at (8.2,0)   {$-\tfrac{1}{n+1}$};
  % cancellation arcs between matching +/- pairs
  \draw[trapred,thick,<->] (b1) to[bend right=45] (a2);
  \draw[trapred,thick,<->] (b2) to[bend right=45] (a3);
  % survivors highlighted
  \node[draw=targetgreen,rounded corners,inner sep=2pt,fit=(a1)] {};
  \node[draw=targetgreen,rounded corners,inner sep=2pt,fit=(bn)] {};
  \node[targetgreen] at (0.0,-0.7) {survives};
  \node[targetgreen] at (8.2,-0.7) {survives};
  \node[trapred]     at (3.4,0.8)  {interior pairs cancel};
\end{tikzpicture}
\end{center}

The closed form $1 - \tfrac{1}{n+1}$ is exact, and it tends to $1$ as
$n \to \infty$. For $n=3$: $\tfrac12 + \tfrac16 + \tfrac{1}{12} = \tfrac{9}{12} =
\tfrac34$, and the formula gives $1 - \tfrac14 = \tfrac34$.
\end{example}

\begin{trapbox}
Telescoping fails the moment a term does \emph{not} have a matching opposite. If
your summand is $a_{k+1} - a_{k-1}$ (a gap of two), the cancellation leaves
\emph{two} terms at each end, not one. Always write out the first three and last
three terms before declaring victory --- the surviving terms are exactly those
that lack a partner.
\end{trapbox}

\section{A field guide to the recurrences in this book}
\label{ch37:recurrences}

Every technique above earns its keep when you unroll a recurrence into a sum. Here
is the workflow, then the catalogue.

\begin{labbox}
\textbf{The unroll-and-sum recipe.}
\begin{enumerate}
  \item Write $T(n)$ in terms of smaller inputs.
  \item Substitute the recurrence into itself a few times until you see the
        pattern of per-level costs.
  \item Express the total as a \emph{sum over levels} (or over $i$).
  \item Identify the sum's shape --- arithmetic, geometric, or harmonic --- and
        apply the matching formula or bound.
\end{enumerate}
Almost every analysis in this book is an instance of this recipe.
\end{labbox}

\subsection{Subtract-a-constant recurrences \texorpdfstring{$\to$}{lead to} arithmetic sums}

A recurrence like $T(n) = T(n-1) + n$ (the cost of inserting into the back of a
sorted list, summed over a build) unrolls into the arithmetic series:
\[
  T(n) = T(0) + \sum_{k=1}^{n} k = \frac{n(n+1)}{2} = \Theta(n^2).
\]
More generally $T(n) = T(n-1) + f(n)$ telescopes to $T(0) + \sum_{k=1}^{n} f(k)$,
and you finish with the integral bound on $\sum f(k)$. Notice this is literally
telescoping (Section~\ref{ch37:telescoping}) followed by a sum.

\subsection{Halve-the-input recurrences \texorpdfstring{$\to$}{lead to} geometric sums}

A recurrence like $T(n) = 2\,T(n/2) + n$ (merge sort) unrolls into a recursion
tree with $\log_2 n$ levels, and the per-level cost is $n$ at \emph{every} level
because the doubling number of subproblems exactly cancels their halving size.
This is the flat case $x = 1$:
\[
  T(n) = \sum_{\text{level } i = 0}^{\log_2 n} n = n\,(\log_2 n + 1) = \Theta(n \log n).
\]
Change the combine cost and the per-level sum becomes a genuine geometric series.
With $T(n) = 2\,T(n/2) + n^2$ the per-level cost \emph{shrinks} geometrically
($x = \tfrac12$): the root's $n^2$ dominates and $T(n) = \Theta(n^2)$. With
$T(n) = 4\,T(n/2) + n$ it \emph{grows} ($x = 2$): the leaves dominate and
$T(n) = \Theta(n^2)$. First-term-wins versus last-term-wins, exactly as in
Section~\ref{ch37:geo-regimes}.

\begin{center}
\begin{tikzpicture}[level distance=11mm,
   level 1/.style={sibling distance=34mm},
   level 2/.style={sibling distance=17mm},
   every node/.style={font=\scriptsize}]
  \node[draw,rounded corners,fill=blue!6] {cost $n$}
    child {node[draw,rounded corners,fill=blue!6] {$n/2$}
      child {node[draw,rounded corners,fill=blue!6] {$n/4$}}
      child {node[draw,rounded corners,fill=blue!6] {$n/4$}}}
    child {node[draw,rounded corners,fill=blue!6] {$n/2$}
      child {node[draw,rounded corners,fill=blue!6] {$n/4$}}
      child {node[draw,rounded corners,fill=blue!6] {$n/4$}}};
  \node[anchor=west,sourceblue] at (3.0,0)    {row sums to $n$};
  \node[anchor=west,sourceblue] at (3.0,-1.1) {row sums to $n$};
  \node[anchor=west,sourceblue] at (3.0,-2.2) {row sums to $n$};
\end{tikzpicture}
\end{center}

The picture makes the flat case obvious: each row of the merge-sort recursion tree
costs $n$, there are about $\log_2 n$ rows, so the sum is the constant $n$ added
$\log_2 n$ times.

\subsection{Random-pivot and tree-build recurrences \texorpdfstring{$\to$}{lead to} harmonic sums}

The expected number of comparisons in randomized quicksort satisfies a recurrence
whose solution is governed by $\sum 1/k$. The clean way to see it: the elements of
rank $i$ and rank $j$ (with $i < j$) are compared with probability
$\tfrac{2}{j-i+1}$, and summing those probabilities over all pairs gives the
expected total. Fix the smaller rank $i$ and let the gap $g = j - i$ run from $1$
up to at most $n-1$; the inner sum over gaps is then a harmonic series:
\[
  \sum_{i<j}\frac{2}{j-i+1}
    \;=\; \sum_{i=1}^{n}\sum_{j=i+1}^{n}\frac{2}{j-i+1}
    \;\le\; \sum_{i=1}^{n} 2\,H_n
    \;=\; 2n\,H_n
    \;=\; \Theta(n \log n),
\]
using $H_n = \Theta(\log n)$ from Section~\ref{ch37:harmonic}. (The middle step
replaces the inner sum $\sum_{g=1}^{n-i}\tfrac{2}{g+1}$, which is at most
$2(H_n - 1) < 2H_n$, by the larger constant $2H_n$ to keep the bookkeeping simple.)
Whenever a recurrence averages over a uniformly random split point, expect a
harmonic series and therefore a $\log n$ factor.

\begin{complexitybox}
The three shapes and their verdicts, all in one place:
\begin{center}
\begin{tabular}{@{}lll@{}}
\toprule
\textbf{Sum} & \textbf{Closed form / bound} & \textbf{Asymptotic} \\
\midrule
$\sum_{k=1}^{n} k$            & $n(n+1)/2$                & $\Theta(n^2)$ \\
$\sum_{k=1}^{n} k^p$          & $\approx n^{p+1}/(p+1)$    & $\Theta(n^{p+1})$ \\
$\sum_{k=0}^{n} x^k,\ x<1$    & $\le 1/(1-x)$             & $\Theta(1)$ (first term wins) \\
$\sum_{k=0}^{n} x^k,\ x=1$    & $n+1$                     & $\Theta(n)$ \\
$\sum_{k=0}^{n} x^k,\ x>1$    & $(x^{n+1}-1)/(x-1)$       & $\Theta(x^n)$ (last term wins) \\
$\sum_{k=1}^{n} 1/k$         & $\ln n + \Theta(1)$        & $\Theta(\log n)$ \\
$\sum_{k=1}^{n} \ln k$       & $n\ln n - n + \Theta(\ln n)$ & $\Theta(n\log n)$ \\
\bottomrule
\end{tabular}
\end{center}
\end{complexitybox}

\begin{historybox}
This appendix corresponds to \textbf{CLRS, Appendix A (``Summations'')}. CLRS
states the same formulas and the integral bounds (their equations for arithmetic,
geometric, and harmonic series, and ``approximation by integrals''), but very
tersely --- it assumes you already believe them. The recursion-tree connections
here preview \textbf{CLRS Chapter 4 (``Divide-and-Conquer'')} and the Master
Theorem. Gauss's fold-in-half trick is a piece of mathematical folklore; the
integral-bound technique is the discrete cousin of the integral test from
calculus, which is exactly how mathematicians decide whether an infinite series
converges.
\end{historybox}

\section{Putting it together: a fully worked analysis}
\label{ch37:worked}

Let us close by analyzing a small algorithm from scratch, using nothing but the
tools above, to show the workflow end to end. Consider:

\begin{algorithm}[H]
\caption{TriangleSums}
\KwIn{An integer $n \ge 1$.}
\KwOut{The value $\sum_{i=1}^{n}\sum_{j=1}^{i} j$.}
$\textit{total} \gets 0$\;
\For{$i \gets 1$ \KwTo $n$}{
  \For{$j \gets 1$ \KwTo $i$}{
    $\textit{total} \gets \textit{total} + j$\;
  }
}
\Return{\textit{total}}\;
\end{algorithm}

\textbf{Counting the work.} The body $\textit{total} \gets \textit{total}+j$ runs
once for each pair $(i,j)$ with $1 \le j \le i \le n$. Those pairs form a triangle:
one cell in row $i=1$, two in row $i=2$, and so on, as the picture shows.

\begin{center}
\begin{tikzpicture}[scale=0.5]
  % a staircase triangle of unit cells: row i has i cells
  \foreach \i in {1,...,6}{
    \foreach \j in {1,...,\i}{
      \draw[fill=structurepurple!18,draw=structurepurple!60!black]
        (\j-1,6-\i) rectangle (\j,7-\i);
    }
    \node[font=\scriptsize] at (-0.7,6.5-\i) {$i=\i$};
    \node[font=\scriptsize,structurepurple!60!black] at (\i+0.6,6.5-\i) {$\i$};
  }
  \node[font=\scriptsize] at (3.0,6.6) {row $i$ has $i$ cells};
\end{tikzpicture}
\end{center}

For a fixed $i$ the inner loop runs $i$ times, so the total number of additions is
the number of cells in the triangle:
\[
  \sum_{i=1}^{n} i = \frac{n(n+1)}{2} = \Theta(n^2),
\]
an arithmetic series (Section~\ref{ch37:arithmetic}). So the running time is
$\Theta(n^2)$.

\textbf{Computing the output.} The returned value is a \emph{nested} sum, and we
peel it from the inside out. The inner sum is itself arithmetic:
\[
  \sum_{j=1}^{i} j = \frac{i(i+1)}{2}.
\]
Substitute and use linearity plus the sum-of-squares formula
(Section~\ref{ch37:squares}):
\[
  \sum_{i=1}^{n}\frac{i(i+1)}{2}
   = \frac12\sum_{i=1}^{n}(i^2 + i)
   = \frac12\!\left(\frac{n(n+1)(2n+1)}{6} + \frac{n(n+1)}{2}\right)
   = \frac{n(n+1)(n+2)}{6}.
\]

\textbf{A sanity check.} For $n = 3$ the algorithm computes
$(1) + (1{+}2) + (1{+}2{+}3) = 1 + 3 + 6 = 10$, and the formula gives
$\tfrac{3\cdot 4\cdot 5}{6} = 10$. The output is $\Theta(n^3)$ while the running
time is only $\Theta(n^2)$ --- a reminder that the \emph{size of a number} and the
\emph{time to compute it} are different sums, and you must do both.

\begin{recoverybox}
If the worked example felt like a lot, that is fine --- it deliberately stacks
three techniques (term-counting, the arithmetic formula, and the sum-of-squares
formula) in one place. In real analysis you will usually need only one of them at
a time. The meta-skill is the recipe: \emph{name the shape, then apply the tool}.
\end{recoverybox}

\section{Exercises}
\label{ch37:exercises}

\begin{exercise}[Term counting warmup]
\label{ch37:ex1}
How many terms are in $\sum_{k=7}^{30} k$? Evaluate the sum two ways: directly
with the arithmetic formula on the shifted range, and as
$\sum_{k=1}^{30} k - \sum_{k=1}^{6} k$. Confirm the answers agree.
\end{exercise}

\begin{exercise}[Constant summands]
\label{ch37:ex2}
Evaluate $\sum_{k=1}^{n} 5$ and $\sum_{k=3}^{n} 5$ as closed forms in $n$. Then
explain in one sentence why $\sum_{k=1}^{n} c \ne c$ for $n > 1$.
\end{exercise}

\begin{exercise}[Linearity practice]
\label{ch37:ex3}
Find a closed form for $\sum_{k=1}^{n}(4k - 3)$. Check your formula at $n=1$ and
$n=4$ by direct addition.
\end{exercise}

\begin{exercise}[Geometric regimes]
\label{ch37:ex4}
Compute $\sum_{k=0}^{10} 3^k$ exactly using the finite geometric formula, and
$\sum_{k=0}^{\infty} (1/3)^k$ exactly using the infinite formula. For each, state
whether the first or the last term dominates.
\end{exercise}

\begin{exercise}[Engineer a telescope]
\label{ch37:ex5}
Show that $\dfrac{1}{k(k+2)} = \dfrac12\left(\dfrac1k - \dfrac{1}{k+2}\right)$, then
evaluate $\sum_{k=1}^{n}\dfrac{1}{k(k+2)}$ in closed form. (Caution: the gap is
two, so \emph{two} terms survive at each end --- write out the first and last few
terms.)
\end{exercise}

\begin{exercise}[Integral bound from scratch]
\label{ch37:ex6}
Using the integral bounds with $f(x) = \sqrt{x}$, prove that
$\sum_{k=1}^{n} \sqrt{k} = \Theta(n^{3/2})$ and identify the leading constant.
\end{exercise}

\begin{exercise}[Harmonic in disguise]
\label{ch37:ex7}
A program does $\sum_{k=1}^{n} \tfrac{n}{k}$ units of work (each pass does $n/k$
work). Show this is $\Theta(n \log n)$, and explain which series from this
appendix you used.
\end{exercise}

\begin{exercise}[Recurrence to sum]
\label{ch37:ex8}
Solve $T(n) = T(n-1) + 2n - 1$ with $T(0) = 0$ by unrolling into a sum. (Hint:
the answer is a familiar perfect square --- recognize the sum of the first $n$
odd numbers.)
\end{exercise}

\begin{exercise}[Which level wins?]
\label{ch37:ex9}
For each recurrence, decide whether the recursion tree's per-level cost shrinks,
stays flat, or grows, and give the resulting $\Theta$: (a) $T(n)=2T(n/2)+1$,
(b) $T(n)=2T(n/2)+n$, (c) $T(n)=3T(n/2)+n$. Tie each answer back to the
geometric-series regime that decides it.
\end{exercise}

\begin{exercise}[Stirling on a walk, written up]
\label{ch37:ex10}
Carry out the Walking-question prompt rigorously: use the integral bounds on
$f(x)=\ln x$ to prove $n\ln n - n < \ln(n!) \le (n+1)\ln(n+1) - n$, and conclude
$\ln(n!) = \Theta(n\log n)$. State one consequence for sorting lower bounds.
\end{exercise}

\begin{exercise}[Challenge: an arithmetic--geometric sum]
\label{ch37:ex11}
Evaluate $\sum_{k=1}^{n} k\,2^k$ in closed form by adapting the
``multiply-by-the-ratio-and-subtract'' telescoping trick from the geometric-series
proof. (Hint: let $S$ be the sum, consider $2S - S$, and you will be left with a
plain geometric series to finish.)
\end{exercise}

\begin{exercise}[Challenge: bounding without a formula]
\label{ch37:ex12}
There is no elementary closed form for $\sum_{k=1}^{n} \tfrac{1}{k^2}$, yet it
\emph{converges}. Use the telescoping bound $\tfrac{1}{k^2} \le
\tfrac{1}{k(k-1)} = \tfrac{1}{k-1} - \tfrac1k$ (for $k \ge 2$) to prove that
$\sum_{k=1}^{n}\tfrac{1}{k^2} < 2$ for every $n$. Why does this confirm the sum is
$\Theta(1)$, unlike the harmonic series?
\end{exercise}
